% =====================================================================
\section{Oxidants and defence}\label{sec:redox}
\figp{fig_redox.png}{Where the peroxide comes from and where it goes. Below: the cycle of the protective enzyme peroxiredoxin, with the side branch that switches it off.}
An oxidant is a molecule that takes electrons from other molecules. Taking an electron from a protein or from DNA damages it. Oxidants are formed unavoidably wherever a cell works with oxygen. The principal one for what follows is hydrogen peroxide \Hp: it is comparatively long-lived and passes through membranes. In small amounts peroxide is an ordinary signal inside the cell; in large amounts it is a poison. \emph{Oxidative stress} is the state in which the influx of oxidants exceeds the capacity to clear them.

The defence has two layers. The first consists of enzymes that decompose peroxide; the fastest of these are the \emph{peroxiredoxins}. The second is a pool of small molecules that donate electrons and thereby return the enzymes to their working state; the principal such pool is \emph{glutathione}.
\facts
\begin{itemize}
\item Sources: the mitochondria (a by-product of energy production)\tU\ and the enzyme Nox2, which produces oxidants deliberately. In neurons Nox2 is switched on by excitation of NMDA receptors; its product leaves the cell and acts on neighbouring ones.\tS
\item The usual concentration of peroxide in a cell is 1--10\,nM.\tS
\item Glutathione is a short molecule of three amino acids; a cell contains five to six orders of magnitude more of it than of peroxide.\tU\ Half the glutathione of the mouse brain turns over in 60--80 hours.\tP\ Its level in the mouse brain varies 1.2-fold over the day.\tS
\item Neurons obtain the raw material for glutathione from \emph{astrocytes}, the brain's supporting cells.\tU
\item In mature neurons the system that switches on protective genes in response to stress itself is much weaker than in astrocytes.\tS
\end{itemize}
\begin{ineq}
The peroxide balance is fast: influx equals efflux within seconds, so the peroxide concentration is an algebraic function of the influx and the state of the defence. The nonlinearity arises because the efflux saturates and, worse, because the defence itself is damaged by excess peroxide (next section).
\end{ineq}
